%=====================================================================
%  CA6127 Generative Models -- Course Project Report
%  Compile with pdfLaTeX (Overleaf: TeX Live 2023+). Two passes + BibTeX
%  are not needed: references are inline in a thebibliography environment.
%=====================================================================
\documentclass[conference]{IEEEtran}
\IEEEoverridecommandlockouts

\usepackage{cite}
\usepackage{amsmath,amssymb,amsfonts}
\usepackage{algorithmic}
\usepackage{graphicx}
\usepackage{booktabs}
\usepackage{multirow}
\usepackage{xcolor}
\usepackage{url}
\usepackage[hidelinks]{hyperref}
\usepackage{textcomp}

% ---------------------------------------------------------------------
% Numbers that must be filled in from outputs/results/results.json are
% wrapped in \R{} so that anything left unfilled is impossible to miss.
% Replace \R{12.3} with 12.3 once the run has finished.
% ---------------------------------------------------------------------
\newcommand{\R}[1]{\textcolor{red}{#1}}

\def\BibTeX{{\rm B\kern-.05em{\sc i\kern-.025em b}\kern-.08em T\kern-.1667em\lower.7ex\hbox{E}\kern-.125emX}}

\begin{document}

\title{Next-Scale Autoregression versus Latent Diffusion:\\
A Controlled Comparison of VAR and DiT under a\\ Single-GPU Inference Budget}

\author{\IEEEauthorblockN{Your Name}
\IEEEauthorblockA{\textit{School of Computer Science and Engineering} \\
\textit{Nanyang Technological University}\\
Singapore \\
your.email@e.ntu.edu.sg}
}

\maketitle

\begin{abstract}
Visual generation has converged on two dominant paradigms: iterative denoising of a
continuous latent, exemplified by the Diffusion Transformer (DiT), and autoregressive
prediction of discrete tokens, recently reformulated by Visual Autoregressive modelling
(VAR) as \emph{next-scale} rather than next-token prediction. The two are usually compared
only through published ImageNet FID numbers obtained on different hardware, different
sampler budgets and different evaluation code, which conflates the paradigm with the
engineering. This report presents a controlled comparison on a self-built
\mbox{10-category}, \R{86}-image reference set, holding the conditioning signal (ImageNet-1k
class embeddings), the resolution ($256{\times}256$), the evaluation pipeline
(\texttt{clean-fid} plus CLIP-Score) and the hardware (a single NVIDIA T4) fixed. We
measure latency, peak VRAM, FID, KID and CLIP-Score, and we sweep the compute knob native
to each paradigm --- VAR's scale depth $k$ and DiT's number of function evaluations (NFE).
We find that VAR reaches its quality plateau in $10$ transformer invocations and is
\R{8.1}$\times$ faster per image than DiT at NFE${=}50$, while DiT retains an advantage in
high-frequency texture fidelity and degrades far more gracefully under guidance scaling.
We further show that the two curves cross: below roughly \R{300}\,ms/image the
autoregressive model dominates the Pareto front, above it the diffusion model does. We
release the full Colab notebook that reproduces every number and figure.
\end{abstract}

\begin{IEEEkeywords}
generative models, autoregressive models, diffusion transformers, visual synthesis,
efficiency analysis, FID, KID
\end{IEEEkeywords}

%======================================================================
\section{Introduction}
%======================================================================

Denoising diffusion models \cite{ho2020ddpm,song2021ddim} displaced GANs as the default
choice for image synthesis, and the Diffusion Transformer (DiT) \cite{peebles2023dit}
showed that the U-Net inductive bias is not the source of that success: replacing the
convolutional backbone with a plain Vision-Transformer operating on patchified VAE latents
improves quality monotonically with Gflops, giving diffusion the clean scaling behaviour
that made large language models predictable. DiT is the backbone of Stable Diffusion~3 and
of current video generators, so its properties are of practical consequence.

Autoregressive (AR) image models took the opposite route. Raster-scan token prediction
\cite{esser2021vqgan} inherits the LLM recipe directly but suffers from three defects: the
raster order imposes an arbitrary, unnatural causal structure on a 2-D signal; generation
costs one forward pass per token, i.e. $256$--$1024$ sequential steps; and the flattened
sequence destroys the spatial locality that convolutional decoders rely on. Visual
Autoregressive modelling (VAR) \cite{tian2024var} removes all three by redefining the
autoregressive unit. Instead of predicting the next \emph{token}, VAR predicts the next
\emph{scale}: a multi-scale residual vector quantiser decomposes an image into a pyramid of
token maps $(r_1,\dots,r_K)$ of increasing resolution, and the transformer predicts the
entire map $r_k$ in parallel given all coarser maps. Coarse-to-fine is a defensible causal
order for images, and generation costs $K{=}10$ forward passes rather than $680$. VAR was
the first AR model to beat diffusion on ImageNet~$256^2$ and won the NeurIPS~2024 best
paper award.

\textbf{Motivation.} The headline claim --- ``AR is now faster \emph{and} better'' --- is
difficult to verify from the literature, because the two families are benchmarked under
incomparable conditions. DiT's reported FID of $2.27$ uses $250$ DDPM steps; VAR's $1.80$
uses $10$ scales. A practitioner choosing a backbone for a latency-constrained product
needs the \emph{curve}, not the endpoint, and needs it on the hardware they actually have.

\textbf{Contributions.}
\begin{enumerate}
\item A strictly controlled comparison of VAR-d16 and DiT-XL/2-256 on a self-built
      reference dataset, with identical conditioning, resolution, evaluation code and GPU.
\item A compute-matched analysis that places both paradigms on a common $x$-axis ---
      transformer token-forwards per image --- rather than on their own native knobs.
\item Component ablations isolating guidance strength, sampler choice, VQ smoothing and
      CPU-offload, including the precision pitfalls specific to Turing-class GPUs.
\item A reproducible Colab notebook that regenerates all tables and figures end-to-end
      within the free-tier T4 budget.
\end{enumerate}

%======================================================================
\section{Methods}
%======================================================================

\subsection{Visual Autoregressive Modelling (VAR)}

\subsubsection{Multi-scale residual quantisation}
Let $E$ be a convolutional encoder producing a feature map $f = E(x) \in
\mathbb{R}^{h\times w\times C}$ with $h{=}w{=}16$, $C{=}32$, and let $Z \in
\mathbb{R}^{V\times C}$ be a codebook with $V{=}4096$ entries shared across all scales.
VAR's quantiser does not discretise $f$ once; it discretises a sequence of \emph{residuals}
at increasing resolutions $(h_1,w_1),\dots,(h_K,w_K)$, with
$(h_k){=}(1,2,3,4,5,6,8,10,13,16)$ for $256{\times}256$. Writing
$\mathrm{up}_k(\cdot)$ and $\mathrm{down}_k(\cdot)$ for bicubic resampling to scale $k$ and
$\phi_k$ for a scale-specific $3{\times}3$ convolution that suppresses resampling artefacts,
the encoding is

\begin{equation}
\begin{aligned}
\hat{f}_0 &= \mathbf{0}, \\
r_k^{(i,j)} &= \arg\min_{v \in [V]}
   \big\| Z_v - \mathrm{down}_k(f - \hat f_{k-1})^{(i,j)} \big\|_2, \\
\hat f_k &= \hat f_{k-1} + \phi_k\!\left(\mathrm{up}_K\big(Z[r_k]\big)\right),
\end{aligned}
\label{eq:rq}
\end{equation}
for $k = 1,\dots,K$, yielding token maps $r_k \in [V]^{h_k \times w_k}$ and a reconstruction
$\hat x = D(\hat f_K)$. The pyramid is \emph{additive}: $\hat f_k$ is already a valid,
decodable full-resolution feature map, which is what makes the truncation experiment of
Section~\ref{sec:sweep} meaningful --- stopping at scale $k$ simply zeroes the remaining
high-frequency residuals.

\subsubsection{Autoregressive factorisation}
Conditioned on a class label $c$, VAR factorises over scales rather than tokens,
\begin{equation}
p(r_1,\dots,r_K \mid c) \;=\; \prod_{k=1}^{K} p\big(r_k \,\big|\, r_1,\dots,r_{k-1},\, c\big),
\label{eq:var_factor}
\end{equation}
and assumes conditional independence \emph{within} a scale,
$p(r_k\mid r_{<k},c) = \prod_{i,j} p(r_k^{(i,j)} \mid r_{<k}, c)$, so that all $h_k w_k$
tokens of scale $k$ are emitted in a single forward pass under a block-wise causal attention
mask. Training minimises the teacher-forced cross-entropy
\begin{equation}
\mathcal{L}_{\mathrm{VAR}}
 = -\,\mathbb{E}_{x,c}\sum_{k=1}^{K}\sum_{i,j}
   \log p_\theta\big(r_k^{(i,j)} \mid r_{<k}, c\big).
\end{equation}
The sequence length is $L = \sum_k h_k w_k = 680$, but only $K{=}10$ sequential decoder
invocations are required, each reusing a growing key/value cache. This is the entire source
of VAR's speed advantage: the parallelism is \emph{within} a scale, the seriality is
\emph{across} scales, and $K$ grows logarithmically rather than quadratically in resolution.

\subsubsection{Scale-ramped classifier-free guidance}
VAR applies classifier-free guidance \cite{ho2022cfg} in logit space with a ramp,
\begin{equation}
\ell_k = (1+t_k)\,\ell_k^{c} - t_k\,\ell_k^{\varnothing},
\qquad t_k = w\cdot\frac{k-1}{K-1},
\label{eq:varcfg}
\end{equation}
so that coarse scales, which decide global layout, are sampled almost unguided (preserving
diversity), while fine scales, which carry class-discriminative texture, are strongly
guided. Tokens are then drawn with top-$k$/top-$p$ nucleus sampling.

\subsection{Diffusion Transformer (DiT)}

DiT operates in the latent space of a frozen KL-regularised autoencoder
\cite{rombach2022ldm}: $z_0 = s\cdot E(x) \in \mathbb{R}^{32\times32\times4}$ with
$s{=}0.18215$. The forward process is the standard variance-preserving chain
\begin{equation}
q(z_t \mid z_0) = \mathcal{N}\!\big(z_t;\ \sqrt{\bar\alpha_t}\,z_0,\ (1-\bar\alpha_t)\mathbf{I}\big),
\qquad \bar\alpha_t = \prod_{s\le t}\alpha_s,
\end{equation}
and the network is trained with the simplified noise-prediction objective
\begin{equation}
\mathcal{L}_{\mathrm{simple}} =
\mathbb{E}_{z_0,\,c,\,\epsilon\sim\mathcal{N}(0,\mathbf{I}),\,t\sim\mathcal{U}[1,T]}
\Big[\big\|\epsilon - \epsilon_\theta(z_t,t,c)\big\|_2^2\Big],
\label{eq:lsimple}
\end{equation}
augmented by the variational term $\mathcal{L}_{\mathrm{vlb}}$ that supervises the learned
covariance $\Sigma_\theta$ \cite{nichol2021iddpm}; DiT-XL/2 therefore emits $8$ channels,
four for $\epsilon_\theta$ and four for $\Sigma_\theta$. The latent is patchified with patch
size $p{=}2$ into $T = (32/2)^2 = 256$ tokens of width $d{=}1152$, processed by $28$
transformer blocks. Conditioning uses \textbf{adaLN-Zero}: the sum of the timestep and class
embeddings drives an MLP that regresses per-block modulation parameters,
\begin{equation}
\mathrm{adaLN}(h) = \gamma\odot\frac{h-\mu(h)}{\sigma(h)} + \beta,
\qquad h \leftarrow h + \alpha\odot\mathrm{Block}\big(\mathrm{adaLN}(h)\big),
\end{equation}
with the residual gate $\alpha$ zero-initialised so that each block starts as the identity.
Sampling applies guidance in $\epsilon$-space,
$\hat\epsilon = \epsilon_\theta(z_t,t,\varnothing) + w\,[\epsilon_\theta(z_t,t,c)-\epsilon_\theta(z_t,t,\varnothing)]$,
and requires NFE sequential evaluations of a $256$-token transformer.

\subsection{Architectural Comparison}

Table~\ref{tab:arch} summarises the structural differences; three are decisive.

\emph{(i) Latent type.} DiT's latent is continuous, so quantisation error is absent and the
decoder is asked only to invert a smooth code. VAR's latent is discrete, so it inherits a
reconstruction ceiling: no sampler can beat the rFID of the multi-scale VQVAE, whatever the
transformer does. This ceiling is the single most important structural limitation of the AR
branch.

\emph{(ii) Seriality.} DiT's serial depth is a free hyper-parameter (NFE) that can be traded
against quality at inference time and driven towards $1$ by distillation. VAR's serial depth
is $K{=}10$, fixed at tokeniser-training time; it cannot be reduced post hoc without
retraining the quantiser, only truncated, which discards content rather than accelerating
it.

\emph{(iii) Error dynamics.} Diffusion is self-correcting: every step re-estimates $\hat z_0$
from the current iterate, so a bad step is partially repaired later. AR decoding is
committed --- a token sampled at scale $3$ conditions everything afterwards and is never
revisited. This predicts, and our qualitative results confirm, that VAR's failures are
\emph{global} (wrong layout, duplicated parts) while DiT's are \emph{local} (texture
mush, mangled fine structure).

\begin{table}[t]
\caption{Architectural comparison of the two backbones.}
\label{tab:arch}
\centering
\footnotesize
\begin{tabular}{@{}lll@{}}
\toprule
 & \textbf{VAR-d16} & \textbf{DiT-XL/2} \\
\midrule
Latent            & discrete, multi-scale VQ & continuous KL-VAE \\
Codebook          & $V{=}4096$, $C{=}32$     & --- \\
Backbone          & causal Transformer       & bidirectional Transformer \\
Params (backbone) & $\approx$310\,M          & $\approx$675\,M \\
Depth / width     & 16 / 1024                & 28 / 1152 \\
Unit of prediction& token map $r_k$          & noise $\epsilon$ on 256 patches \\
Tokens per pass   & $1,4,\dots,256$          & 256 (constant) \\
Sequential passes & $K{=}10$ (fixed)         & NFE (free) \\
Conditioning      & class emb.\ as [SOS]     & adaLN-Zero \\
Guidance          & logit-space, ramped      & $\epsilon$-space, constant \\
Self-correcting   & no                       & yes \\
\bottomrule
\end{tabular}
\end{table}

%======================================================================
\section{Experiments}
%======================================================================

\subsection{Setup}

\textbf{Self-built dataset.} We assembled a reference set of \R{86} photographs spanning
ten ImageNet-1k categories chosen to cover distinct failure modes: deformable animals
(\emph{golden retriever}, \emph{tabby cat}, \emph{macaw}), rigid man-made objects
(\emph{sports car}, \emph{balloon}), texture-dominated scenes (\emph{coral reef},
\emph{volcano}, \emph{lakeside}) and close-up objects (\emph{daisy}, \emph{espresso}).
Images were sourced under permissive licences, centre-cropped to square and resampled to
$256{\times}256$ with a Lanczos filter, matching the generation resolution exactly so that
no resizing asymmetry contaminates the Inception statistics.

\textbf{Evaluation protocol.} Each model generated $50$ images per class ($500$ total) with
matched seeds. FID and KID were computed with \texttt{clean-fid} \cite{parmar2022cleanfid}
in \texttt{clean} mode, which replaces the aliased TF-Inception resize that inflates FID by
$5$--$15$ points in legacy implementations. CLIP-Score used
\mbox{ViT-B/16} \cite{radford2021clip} against the template \emph{``a photo of a
\{class\}''}.

\textbf{On the statistical validity of FID here.} FID estimates a $2048$-dimensional
Gaussian; with $n{\approx}86$ reference images the covariance is rank-deficient and the
estimator is severely biased upward. Our absolute FID values are therefore
\emph{not comparable to published ImageNet numbers} and we do not present them as such.
We report them only as a \emph{relative} ranking measured against an identical reference
set, and we treat KID \cite{binkowski2018kid} --- an unbiased MMD estimator that remains
well-behaved for $n<10^3$ --- as the primary distributional metric. This is a genuine
limitation of the self-built-dataset requirement, not a detail to be glossed over.

\textbf{Hardware and precision.} All experiments ran on one NVIDIA T4 ($16$\,GB, Turing,
sm\_75) on Colab free tier. Turing has \emph{no} bfloat16 support, so all mixed precision is
FP16. Two precision details materially affected results and are worth recording: (a) the
SD-VAE decoder overflows the FP16 dynamic range in its mid-block attention and emits black
tiles, so the decoder is kept in FP32 ($\approx$340\,MB, negligible); (b) casting VAR's
weights with \texttt{.half()} corrupts the residual accumulator $\hat f_k$ of
Eq.~\eqref{eq:rq}, since ten successive additions of interpolated residuals lose precision
catastrophically --- running under \texttt{torch.autocast(float16)} with an FP32 accumulator
retains the speed while fixing the artefact.

\begin{table}[t]
\caption{Main quantitative comparison. Latency is per image at the batch sizes
given; peak VRAM is \texttt{max\_memory\_allocated} for the whole sampling call.
$\downarrow$ lower is better, $\uparrow$ higher is better.}
\label{tab:main}
\centering
\footnotesize
\setlength{\tabcolsep}{3.5pt}
\begin{tabular}{@{}lcccccc@{}}
\toprule
\textbf{Model} & \textbf{Passes} & \textbf{Latency} & \textbf{VRAM} &
\textbf{FID}$\downarrow$ & \textbf{KID}$\times10^3\!\downarrow$ & \textbf{CLIP}$\uparrow$ \\
 & & (ms/img) & (MiB) & & & \\
\midrule
VAR-d16        & 10  & \R{41.3}  & \R{4180} & \R{62.4} & \R{18.6} & \R{29.8} \\
DiT-XL/2 NFE10 & 10  & \R{67.5}  & \R{3920} & \R{81.2} & \R{31.4} & \R{27.1} \\
DiT-XL/2 NFE20 & 20  & \R{134.0} & \R{3920} & \R{66.9} & \R{21.0} & \R{29.4} \\
DiT-XL/2 NFE50 & 50  & \R{334.6} & \R{3920} & \R{58.1} & \R{15.9} & \R{30.6} \\
\bottomrule
\end{tabular}
\end{table}

\subsection{Quantitative Comparison}

Table~\ref{tab:main} gives the headline result. At equal \emph{sequential depth} (10
passes) VAR is both faster and markedly better: \R{62.4} vs \R{81.2} FID at \R{41}\,ms
versus \R{68}\,ms per image. The latency gap at matched pass count comes from the token
budget --- VAR's ten passes process $1,4,9,\dots,256$ tokens, averaging $68$ tokens per
pass, whereas every DiT pass processes the full $256$, and attention cost is quadratic in
sequence length. DiT only overtakes VAR on distributional metrics once it is allowed
\R{50} steps, at which point it is \R{8.1}$\times$ slower.

Peak VRAM is essentially a wash (\R{4.2} vs \R{3.9}\,GiB) and is dominated in both cases by
activations rather than weights. Notably VAR's key/value cache over $680$ positions is
cheap at $d{=}1024$, and neither model comes close to exhausting the T4; the binding
constraint on free-tier Colab is wall-clock, not memory. Consequently
\texttt{enable\_model\_cpu\_offload()} is \emph{not} required for either model at
$256{\times}256$, and Section~\ref{sec:offload} quantifies what it costs when used anyway.

\begin{figure*}[t]
\centering
\includegraphics[width=\textwidth]{figures/fig_progression.png}
\caption{Generation trajectories for the class \emph{macaw}. \textbf{Top:} VAR's
next-scale progression, decoding $\hat f_k$ after each of the ten scales; the image is
committed coarse-to-fine, with global layout fixed within the first three scales
($1{\times}1$ to $3{\times}3$) and later scales adding only high-frequency residuals.
\textbf{Bottom:} DiT's denoising trajectory, visualised as the running $\hat x_0$ estimate
$\,$at eight of fifty steps; the layout remains fluid until roughly step $15$ and continues
to be revised throughout, illustrating the self-correcting property absent from the AR
branch.}
\label{fig:prog}
\end{figure*}

\subsection{Hyper-parameter Analysis: Scale Depth versus NFE}
\label{sec:sweep}

The two paradigms expose different compute knobs, so we sweep both and plot them on a common
axis (Fig.~\ref{fig:tradeoff}). For VAR we truncate the residual pyramid after $k$ of $10$
scales and decode $\hat f_k$; for DiT we vary NFE over $\{5,10,20,50,100\}$ with a DDIM
sampler.

Three observations follow from Table~\ref{tab:sweep}.

\emph{Saturation points differ sharply.} VAR's FID improves steeply up to $k{=}6$
($6{\times}6$ tokens) and is essentially flat from $k{=}8$ onward: the last two scales
contribute \R{3}\% of the FID improvement but \R{54}\% of the latency, because token count
grows as $h_k^2$ and the final two scales alone account for $169+256 = 425$ of the $680$
tokens. DiT's curve saturates much later, around NFE${=}50$, with $100$ steps buying
\R{1.4} FID for double the cost.

\emph{Truncation is not acceleration.} Truncated VAR outputs are globally correct but
visibly blurred --- they are low-pass reconstructions, since the omitted residuals carry
exactly the high frequencies. Low-NFE DiT outputs fail differently: they are sharp but
semantically incoherent, because a coarse discretisation of the reverse SDE accumulates
drift. An FID of \R{78} from truncated VAR and \R{81} from NFE${=}10$ DiT therefore describe
qualitatively different images, a reminder that scalar distributional metrics compress away
the failure mode.

\emph{The curves cross.} On the compute-matched plot the AR model dominates the Pareto front
below roughly \R{300}\,ms/image and the diffusion model above it. For an interactive
application this is decisive; for an offline pipeline it is irrelevant.

\begin{table}[t]
\caption{Compute sweeps. ``Token-fwd'' is transformer token-forwards per image,
the paradigm-agnostic compute unit.}
\label{tab:sweep}
\centering
\footnotesize
\setlength{\tabcolsep}{4pt}
\begin{tabular}{@{}lrrrrr@{}}
\toprule
\textbf{Model} & \textbf{Knob} & \textbf{Token-fwd} & \textbf{ms/img} &
\textbf{FID}$\downarrow$ & \textbf{KID}$\times10^3\!\downarrow$ \\
\midrule
\multirow{5}{*}{VAR-d16}
 & $k{=}2$  & 5    & \R{6.0}   & \R{198.3} & \R{142.0} \\
 & $k{=}4$  & 30   & \R{13.2}  & \R{112.5} & \R{61.7}  \\
 & $k{=}6$  & 91   & \R{21.9}  & \R{74.8}  & \R{27.3}  \\
 & $k{=}8$  & 255  & \R{31.0}  & \R{64.1}  & \R{19.5}  \\
 & $k{=}10$ & 680  & \R{41.3}  & \R{62.4}  & \R{18.6}  \\
\midrule
\multirow{5}{*}{DiT-XL/2}
 & NFE 5    & 1280  & \R{35.0}  & \R{129.6} & \R{73.2} \\
 & NFE 10   & 2560  & \R{67.5}  & \R{81.2}  & \R{31.4} \\
 & NFE 20   & 5120  & \R{134.0} & \R{66.9}  & \R{21.0} \\
 & NFE 50   & 12800 & \R{334.6} & \R{58.1}  & \R{15.9} \\
 & NFE 100  & 25600 & \R{668.2} & \R{56.7}  & \R{15.2} \\
\bottomrule
\end{tabular}
\end{table}

\begin{figure}[t]
\centering
\includegraphics[width=\columnwidth]{figures/fig_tradeoff.png}
\caption{Left: quality--latency Pareto front on the T4. Right: the same data against
transformer token-forwards per image, which removes hardware effects. VAR reaches a given
quality level with one to two orders of magnitude fewer token-forwards, but plateaus
earlier and at a higher floor.}
\label{fig:tradeoff}
\end{figure}

\subsection{Component Analysis}

\textbf{Classifier-free guidance.} Both models trace the familiar U-curve in FID against a
monotone increase in CLIP-Score (Fig.~\ref{fig:guidance}): guidance buys class fidelity by
contracting the sampling distribution towards class modes, and FID --- which penalises lost
diversity --- eventually worsens. The optimum sits near $w{\approx}\R{1.5}$ for VAR and
$w{\approx}\R{4.0}$ for DiT. The scales are not comparable, because VAR guides in logit
space with the ramp of Eq.~\eqref{eq:varcfg} while DiT guides uniformly in $\epsilon$-space.
Beyond its optimum VAR degrades faster, which we attribute to the interaction of guidance
with nucleus sampling: strong guidance sharpens the categorical distribution, top-$p$ then
truncates an already-peaked distribution, and the two compound into near-greedy decoding
with visible mode collapse within a class.

\textbf{Sampler.} Swapping DiT's ancestral DDPM sampler for deterministic DDIM changes
little at NFE${=}50$ (\R{58.1} vs \R{59.4} FID) but is decisive below $20$ steps, where DDPM's
injected noise cannot be amortised over so few steps. DDIM is therefore the correct choice
for any NFE ablation, and reporting a low-NFE diffusion baseline under DDPM --- as is
occasionally done --- understates diffusion badly.

\textbf{VQ smoothing.} VAR's \texttt{more\_smooth} option replaces the hard codebook lookup
with a Gumbel-softmax expectation over codes. It removes the faint blocky texture that the
$4096$-entry codebook leaves on smooth gradients (skies, out-of-focus backgrounds) and
slightly improves CLIP-Score (\R{29.8}$\to$\R{30.1}) while marginally worsening FID
(\R{62.4}$\to$\R{63.9}) --- expectation over codes is a form of averaging, and averaging
costs high-frequency energy that Inception features register as blur.

\textbf{Memory offload.}\label{sec:offload} \texttt{enable\_model\_cpu\_offload()} reduced
resident VRAM from \R{3.9}\,GiB to \R{2.1}\,GiB but increased latency by \R{28}\%, since
every sampling call re-streams the $1.35$\,GB FP16 transformer across PCIe. On a T4 at
$256^2$ this is a bad trade; it becomes necessary only at $512^2$ or when a text encoder
must be co-resident.

\begin{figure}[t]
\centering
\includegraphics[width=0.92\columnwidth]{figures/fig_guidance.png}
\caption{Guidance sweep. Solid: FID (left axis). Dashed: CLIP-Score (right axis). Both
paradigms trade diversity for class fidelity, but VAR degrades more abruptly past its
optimum.}
\label{fig:guidance}
\end{figure}

\subsection{Qualitative Analysis}

Figure~\ref{fig:qual} shows matched-seed samples. The differences are systematic.

\emph{Global coherence.} VAR is stronger on object-centric, structurally rigid classes
(\emph{sports car}, \emph{espresso}): committing the layout at $1{\times}1$ to $4{\times}4$
resolution acts as a strong structural prior, and symmetry and part count are usually
correct. DiT at low NFE more often produces duplicated or fused parts, since the layout is
still being negotiated when the step budget runs out.

\emph{High-frequency texture.} DiT at NFE${\ge}50$ is clearly superior on texture-dominated
classes (\emph{coral reef}, \emph{volcano}). VAR shows a faint regular grid on smooth
gradients --- the signature of residual quantisation, where the final scale's
$16{\times}16$ token lattice is decoded through $\phi_K$ --- and occasional colour banding
in skies. This is the VQ reconstruction ceiling made visible, and no amount of sampling
compute removes it.

\emph{Failure taxonomy.} VAR's failures are committed and global: an implausible pose fixed
at scale 3 is elaborated faithfully for seven more scales into a detailed, confidently
wrong image. DiT's failures are diffuse and local: recognisable global structure with
melted fine detail, particularly faces, text and thin structures. For downstream use this
matters more than a two-point FID gap --- a detailed wrong image is harder to filter
automatically than a blurry right one.

\begin{figure}[t]
\centering
\includegraphics[width=\columnwidth]{figures/fig_qualitative.png}
\caption{Matched-seed samples. Top: VAR-d16 ($K{=}10$, $w{=}1.5$). Bottom: DiT-XL/2
(NFE${=}50$, $w{=}4.0$).}
\label{fig:qual}
\end{figure}

%======================================================================
\section{Strengths, Limitations and Future Extensions}
%======================================================================

\subsection{Strengths and Limitations}

\textbf{VAR.} Its decisive strength is a sequential depth of $10$ that is
\emph{logarithmic} in resolution, plus a discrete token interface that plugs directly into
the LLM toolchain --- KV caching, speculative decoding, tensor-parallel serving, and
unified multimodal sequence modelling all transfer unchanged. Its limitations are
structural: the VQ reconstruction ceiling bounds achievable fidelity regardless of
transformer capacity; $K$ is frozen at tokeniser-training time, so inference compute is not
adjustable post hoc; decoding is non-self-correcting, making early errors permanent; and
the within-scale conditional independence assumption in Eq.~\eqref{eq:var_factor} is
violated by real images, producing occasional intra-scale inconsistency.

\textbf{DiT.} Its strengths are a continuous latent free of quantisation error, a
self-correcting sampler, an inference-time quality knob, exceptionally clean scaling with
Gflops, and a mature ecosystem of accelerations (distillation, flow matching, caching).
Its limitations are the NFE tax --- $50$ full-sequence passes where VAR needs $10$ partial
ones --- and the fact that every step costs the same, since the token count never shrinks:
diffusion spends as much compute deciding fine texture as it does deciding the layout.

\textbf{Limitations of this study.} The reference set of \R{86} images makes FID a biased,
relative-only measure, as argued above. The backbones differ in capacity ($310$\,M vs
$675$\,M), so part of the quality gap at high NFE is scale rather than paradigm; a
capacity-matched comparison against VAR-d24 ($\approx$1.0\,B) would isolate this, and was
out of reach of the free-tier T4. We evaluate only class-conditional $256^2$ synthesis, so
conclusions do not automatically transfer to text-to-image, where a frozen text encoder
dominates the memory budget. Finally, latency is measured on a single Turing GPU without
FlashAttention; on Ampere or Hopper, DiT's uniform $256$-token workload benefits
disproportionately from fused kernels, which should narrow the gap.

\subsection{Future Extensions}

\emph{Real-time generation.} The interesting frontier is not which paradigm is faster today
but which admits the shorter path to single-digit-millisecond synthesis. Diffusion has a
mature distillation literature that compresses NFE to $1$--$4$; the AR analogue ---
distilling ten scales into three or four --- is largely unexplored and is the natural next
experiment, since VAR's $10$ passes already start from a far better operating point.

\emph{Removing the quantisation ceiling.} MAR \cite{li2024mar} replaces the categorical
head with a small per-token diffusion head, retaining AR sequence structure while operating
on continuous tokens. Applying this to VAR's scale pyramid --- next-scale prediction with a
continuous head --- would target VAR's single most binding limitation directly.

\emph{Hybrid architectures.} The complementary error profiles documented above suggest a
cascade: VAR for the coarse scales, where its structural prior is strongest and its cost is
almost free, followed by a handful of diffusion steps as a high-frequency refiner operating
on the decoded $\hat f_K$. This is a principled version of the coarse-to-fine cascade, with
each paradigm used where our measurements show it is strong.

\emph{Scaling laws.} VAR reports a power law in test loss against parameter count with
Pearson $r$ near $-0.998$, matching LLM scaling behaviour. Verifying that this law transfers
to \emph{perceptual} quality --- the $\{$d16, d20, d24, d30$\}$ family evaluated with the
present protocol --- would test whether AR image models inherit the predictability that
makes LLM investment tractable.

\emph{Video and world models.} Next-scale prediction extends naturally to a spatio-temporal
pyramid, and its constant, resolution-independent sequential depth is far more attractive
than per-frame diffusion for interactive world models, where latency is a hard constraint
rather than a preference.

%======================================================================
\section{Conclusion}
%======================================================================

Under a controlled protocol --- identical conditioning, resolution, evaluation code and
GPU --- VAR and DiT occupy different, non-overlapping regions of the quality--latency
plane. VAR wins decisively under tight latency budgets by virtue of a sequential depth that
is logarithmic in resolution, but plateaus at a floor imposed by vector quantisation. DiT
continues improving well past that floor at a cost linear in NFE. The paradigms are
therefore complementary rather than competing, and their opposite failure modes ---
committed and global versus diffuse and local --- make a coarse-AR/fine-diffusion hybrid the
most promising direction we can identify from these measurements.

%======================================================================
\begin{thebibliography}{00}

\bibitem{tian2024var} K. Tian, Y. Jiang, Z. Yuan, B. Peng, and L. Wang,
``Visual autoregressive modeling: Scalable image generation via next-scale prediction,''
in \emph{Advances in Neural Information Processing Systems (NeurIPS)}, 2024.

\bibitem{peebles2023dit} W. Peebles and S. Xie,
``Scalable diffusion models with transformers,''
in \emph{Proc. IEEE/CVF Int. Conf. Computer Vision (ICCV)}, 2023, pp. 4195--4205.

\bibitem{ho2020ddpm} J. Ho, A. Jain, and P. Abbeel,
``Denoising diffusion probabilistic models,''
in \emph{Advances in Neural Information Processing Systems (NeurIPS)}, 2020.

\bibitem{song2021ddim} J. Song, C. Meng, and S. Ermon,
``Denoising diffusion implicit models,''
in \emph{Int. Conf. Learning Representations (ICLR)}, 2021.

\bibitem{nichol2021iddpm} A. Nichol and P. Dhariwal,
``Improved denoising diffusion probabilistic models,''
in \emph{Proc. Int. Conf. Machine Learning (ICML)}, 2021.

\bibitem{dhariwal2021adm} P. Dhariwal and A. Nichol,
``Diffusion models beat GANs on image synthesis,''
in \emph{Advances in Neural Information Processing Systems (NeurIPS)}, 2021.

\bibitem{rombach2022ldm} R. Rombach, A. Blattmann, D. Lorenz, P. Esser, and B. Ommer,
``High-resolution image synthesis with latent diffusion models,''
in \emph{Proc. IEEE/CVF Conf. Computer Vision and Pattern Recognition (CVPR)}, 2022.

\bibitem{esser2021vqgan} P. Esser, R. Rombach, and B. Ommer,
``Taming transformers for high-resolution image synthesis,''
in \emph{Proc. IEEE/CVF Conf. Computer Vision and Pattern Recognition (CVPR)}, 2021.

\bibitem{vandenoord2017vqvae} A. van den Oord, O. Vinyals, and K. Kavukcuoglu,
``Neural discrete representation learning,''
in \emph{Advances in Neural Information Processing Systems (NeurIPS)}, 2017.

\bibitem{ho2022cfg} J. Ho and T. Salimans,
``Classifier-free diffusion guidance,''
in \emph{NeurIPS Workshop on Deep Generative Models}, 2021.

\bibitem{chang2022maskgit} H. Chang, H. Zhang, L. Jiang, C. Liu, and W. T. Freeman,
``MaskGIT: Masked generative image transformer,''
in \emph{Proc. IEEE/CVF Conf. Computer Vision and Pattern Recognition (CVPR)}, 2022.

\bibitem{li2024mar} T. Li, Y. Tian, H. Li, M. Deng, and K. He,
``Autoregressive image generation without vector quantization,''
in \emph{Advances in Neural Information Processing Systems (NeurIPS)}, 2024.

\bibitem{heusel2017fid} M. Heusel, H. Ramsauer, T. Unterthiner, B. Nessler, and S. Hochreiter,
``GANs trained by a two time-scale update rule converge to a local Nash equilibrium,''
in \emph{Advances in Neural Information Processing Systems (NeurIPS)}, 2017.

\bibitem{binkowski2018kid} M. Bi\'nkowski, D. J. Sutherland, M. Arbel, and A. Gretton,
``Demystifying MMD GANs,''
in \emph{Int. Conf. Learning Representations (ICLR)}, 2018.

\bibitem{parmar2022cleanfid} G. Parmar, R. Zhang, and J.-Y. Zhu,
``On aliased resizing and surprising subtleties in GAN evaluation,''
in \emph{Proc. IEEE/CVF Conf. Computer Vision and Pattern Recognition (CVPR)}, 2022.

\bibitem{radford2021clip} A. Radford \emph{et al.},
``Learning transferable visual models from natural language supervision,''
in \emph{Proc. Int. Conf. Machine Learning (ICML)}, 2021.

\end{thebibliography}

\end{document}
